% 對應英文：sections/01_abstract.tex
% TODO: 翻譯後請同步更新英文版的修改

預測下週哪些加密貨幣將超越同類資產，本質上是一個橫截面排序問題，
然而過去的研究普遍採用不適合排序目標的時間序列損失函數。
本文證明，這種目標函數與評估標準之間的不匹配，
是深度學習模型在此類任務中表現遜於普通最小二乘法（OLS）的主要原因。

本文的核心方法貢獻，在於將 ListNet 排序損失的目標訊號從原始報酬替換為
橫截面排名百分位：每個資產依其在當週橫截面中的排名，被賦予介於 $[0, 1]$
的監督信號（Poh 等人，2021），無論當週報酬的絕對波動幅度為何，
每個資產對損失梯度的貢獻均等。

本文在一個涵蓋 324 週 $\times$ 86 種加密貨幣的週資料面板上，
評估三種深度學習架構——具有跨資產注意力的 LSTM、採用完整變數選擇網路的時序融合轉換器（TFT），
以及全新設計的橫截面門控網路（CS-Gated）——並與六種傳統基準模型進行比較。

最佳模型 CS-Gated（以鏈上特徵訓練）在年化夏普比率上達到 2.75（等權多空組合），
而 LSTM 搭配川普社群媒體與 ETF 資金流特徵則達到 2.13。
整體結果顯示，深度學習在正確的排序損失下可與傳統方法競爭，
但不同模型對不同訊號型態各有優勢，隨機森林在高維資訊集下仍具強勁表現。
消融研究（ablation study）確認，排名正規化搭配梯度累積，
是促成績效改善的關鍵機制。
